\section{Marked fields and spatial aggregation}\label{supp:sec:fields}
The distinction between a labelled field and its aggregate is already visible
in a finite family. This appendix gives the finite inequalities first; all
passages to countably many marks are then controlled by a single tail event.

\subsection{Uniform and averaged dictionary comparisons}
The marginal-ceiling calculation in the proof of
\cref{thm:word-field} gives the finite inequality
\[
 \EE\Delta_{\mathcal W}\ll
 a\bigl(\mathcal M_B^{\rm val}+|D_Y^{\rm val}|\bigr)
 +a^2\bigl(NB+E_Y+R_2^{\mathrm{val},[1/a]}(N,B)\bigr)
 +\Lambda\Omega(\mathcal W),
\]
where $a=m2^{-B}$ and $\Lambda=Na$.
The comparison retains the full site--word index set; only the error
terms have been summed over words. We verify here the substitutions
and the uniform absorption of secondary terms.
With $T=N/\Lambda$ and $Q_B=Nm/\Lambda\ge T$, the three capped-profile
terms become, respectively,
\[
 N^{-1/3+\varepsilon}\Lambda^{11/6}m^{1/6},\qquad
 N^{-1/3+\varepsilon}\Lambda^{4/3}m^{2/3},\qquad
 N^{-1/3+\varepsilon}\Lambda.
\]
These substitutions are exact because $a=\Lambda/N$ and
$Q_B=Nm/\Lambda$. The restrictions $1\le m\le2^B$ also give
$0<a\le1$, so the arithmetic cap $T=1/a$ is always at least one.
The secondary edge contribution is absorbed in every admissible
intensity regime, not only at criticality: $m\ge\Lambda/N$ implies
\[
 \frac{\Lambda^2}{N}
 \le N^{-1/3}\Lambda^{4/3}m^{2/3}.
\]
Factors $N^{o(1)}$ in these smaller terms are absorbed by fixing a
smaller arithmetic error exponent first.
At bounded positive $\Lambda$ and $m\le N^{1/2-\eta}$, the powers of
$N$ in the three terms, with $\varepsilon=\eta/6$, are at most
\[
 -\frac14,\qquad -\frac\eta2,\qquad
 -\frac13+\frac\eta6\le-\frac\eta2\quad(0<\eta<1/2).
\]
Thus all three are $O(N^{-\eta/2})$, uniformly in the dictionary. No independent random marks
are added to the multiplicative sequence, and no word identity is lost
when the error terms are summed.

\subsubsection*{Averaging the finite dictionary comparison}
\label{supp:sec:typical-details}
We verify \eqref{eq:typical-finite} before using any asymptotic profile.
Fix $N\ge2$, $B\ge1$, $Y>2B$, a deterministic start mask
$A\subset I_N$, and a random size-$m$ dictionary, $1\le m\le2^B$,
independent of the prime signs. Only the inclusion probabilities
\eqref{eq:dictionary-selection} are needed; both ensembles in the article
have them. For each realization of the dictionary, apply the process
inequality on the finite coordinate set $A\times\mathcal W$.

\paragraph{Deletion and the exact graph.}
Let $D=D_Y^{\rm val}\cap A$. Removing real indicators has mean cost,
after averaging the dictionary and the prime environment,
\[
 \sum_{x\in D}\EE_f\PP_{\mathcal W}(\mathbf f_x\in\mathcal W)=a|D|.
\]
The Poisson target has total deleted mean $a|D|$. This uses no uniformity
of a bad window's law. At a retained site every coordinate has conditional
mean $p_B=2^{-B}$. Its neighbours include all words at its own site,
at each overlapping support and at each support sharing an odd prime
coordinate above $Y$. Families with no edge depend on disjoint remaining
prime signs, so their conditional independence is exact. The site graph
and its good set do not depend on $\mathcal W$.
The first Chen--Stein sum, including each diagonal coordinate, is bounded
by $Ca^2(|A|B+E_Y(A))$.

\paragraph{The second sum.}
At the same site its off-diagonal terms vanish by mutual exclusivity;
a coordinate's own diagonal belongs to the first sum only. For distinct
sites the sum over word labels is exactly
$\PP_f(\mathbf f_x,\mathbf f_y\in\mathcal W\mid\mathcal F_Y)$.
After both averages it is $b+(a-b)\PP_f(\mathbf f_x=\mathbf f_y)$, as in
\eqref{eq:selected-pair}. At a good local pair with displacement $d<B$,
write the union values as $z_0,\ldots,z_{B+d-1}$. The equations
$z_j=z_{j+d}$, $0\le j<B$, form the $d$ paths with residue classes
modulo $d$. There are $B+d$ vertices and $B$ edges. A prime above
$Y>2B$ can occur at only one union vertex, and each vertex has such
an odd coordinate, so the conditional vector is uniform. Its collision
probability is exactly $2^{-B}$ on every fibre.

For disjoint supports, let $v(n)$ denote the complete binary valuation
vector at the finite set of primes dividing the union. The equality
matrix has rows $v(x-1+j)+v(y-1+j)$. Its coefficient kernel maps
injectively by
\[
 t\longmapsto(t,t),\qquad
 \sum_jt_j\bigl(v(x-1+j)+v(y-1+j)\bigr)=0,
\]
into the coefficient kernel of the stacked value matrix. The equality
right-hand side is zero, so its probability is exactly
$2^{-B+\kappa}$, regardless of whether either value window is individually
uniform. This proves \eqref{eq:word-collision} with the full, not the
large-prime-truncated, nullity. It is used only after averaging the prime
fibres. Consequently a separated edge costs at most
\[
 a^2+ap_B+ap_B(2^{\rho_B^{\rm val}}-1)
 \le 2a^2+ap_B(2^{\rho_B^{\rm val}}-1).
\]
The local sum is $O(a^2|A|B)$ and the separated sum is
$O(a^2E_Y(A)+ap_BR_2^{\rm val}(A,B))$. All site pairs here are ordered;
the two orientations are included. Adding the two deletion costs proves
\eqref{eq:typical-finite} with an absolute implied constant. This finite
inequality also covers $m=1$, $B=1$ and the empty mask.

\paragraph{Normalization and the order of averaging.}
With $a=\Lambda/N$ and $p_B=\Lambda/(Nm)$, the two terms of the coarse
profile give, separately,
\[
 ap_BN^{5/3}=\frac{\Lambda^2}{m}N^{-1/3},\qquad
 ap_BN^{2/3}2^B=\Lambda N^{-1/3}.
\]
The factor $N^\varepsilon$ is inserted only after these exact identities.
For an independent-letter word field, distinct disjoint supports need no
edge, and every local collision again has probability $p_B$. Its averaged
Poisson error is therefore $O(a^2|A|B)$, without a restriction on $m$.
This is a fresh finite calculation, not an application of the uniform
dictionary theorem outside its range.

The supremum over target tests is taken for each fixed dictionary and
prime fibre \emph{before} averaging their distances. Finite sums over the
auxiliary dictionary law then give the claimed mean bound. The class of
good dictionaries supplied by Markov controls all coordinate restrictions
simultaneously by contraction. No mask may use the observed prime signs.
A dictionary-dependent deterministic restriction is harmless once that
dictionary has been selected. No coupling over different scales, or
almost-sure choice of a good dictionary at every scale, is inferred.

\subsection{Local pairs and the role of the maximal support}\label{supp:sec:local-marks}
From here on, $p=2^{-L}$, $q_e=2^{-e-1}$, $p_e=pq_e$,
$\lambda=Np$ and $0\le e\le E$. The largest row length is
$Q=L+E+1$; its support has $Q+1$ vertices. An unsigned mark of excess $e$
imposes $L+e+1$ relative equations; prescribing its sign fixes all
$L+e+2$ values on the same support.

A good maximal support contains no $Y$-defective vertex and has diameter
$Q<Y$. Every exact-mark system on that support consequently has full rank
above $Y$, with conditional marginal $p_e$. If two actual supports intersect,
their union has diameter at most $2Q<Y$; every vertex of the union then has a
prime coordinate private to that union. For two compatible exact runs the
only possible intersecting cases are a common terminal/border vertex or a
common terminal edge. Equivalently, for $x<y$, they satisfy
$y-x=L+e+1$ or $y-x=L+e$. The union of the two constraint trees has cycle
rank at most one, so its probability is at most $2p_ep_f$.

The actual supports can also be disjoint while the \emph{maximal} supports
intersect. In that case all their private coordinates are still distinct,
and the probability is $p_ep_f$. Thus the bound
\begin{equation}\label{supp:eq:local-marked}
 \PP(K_{x,e}K_{y,f}=1\mid\mathcal F_Y)\le 2p_ep_f
 \qquad(0<|x-y|\le Q)
\end{equation}
holds for every good pair; it does not assert that only two separations are
possible among pairs of intersecting maximal supports. At the same site,
different exact marks are mutually exclusive.

For $|x-y|>Q$, zero-extension of the coefficient vectors embeds the mixed
relation space into the two-window space at row length $Q$. Therefore
\begin{equation}\label{supp:eq:mixed-domination}
 \PP(K_{x,e}K_{y,f}=1)\le p_ep_f\,2^{\rho_Q(x,y)}.
\end{equation}
This is an averaged, not necessarily fibrewise, estimate: the small-prime
phase in the affine formula can vary with the environment.


For signed marks, put $p_{e,s}=p_e/2$ and let $a=L+e$, $b=L+f$.
On a good local union the raw value vector is uniform. For $x<y$,
$d=y-x$, the exact joint probabilities are
\[
 \PP(K^s_{x,e}K^t_{y,f}=1\mid\mathcal F_Y)=
 \begin{cases}
  0,&d<a,\\
  4p_{e,s}p_{f,t}\1_{\{t=-s\}},&d=a,\\
  2p_{e,s}p_{f,t}\1_{\{t=s\}},&d=a+1,\\
  p_{e,s}p_{f,t},&a+1<d\le Q.
 \end{cases}
\]
The final case concerns disjoint actual supports inside intersecting
maximal supports. The factor four, rather than two, is needed when signs
are prescribed. After summing signs one recovers
\eqref{supp:eq:local-marked}. For separated signed systems, zero-extension
instead gives the bound
$p_{e,s}p_{f,t}2^{\rho^{\rm val}_{Q+1}(x,y)}$.
\Cref{prop:absolute-profile} controls its weighted sum.

\subsection{The directional multivariate Stein factor}\label{supp:sec:stein-factor}
Let $\pi=\bigotimes_{i=1}^d\Pois(\lambda_i)$, with $d\ge2$,
$\lambda_i>0$, $\lambda'=\sum_i\lambda_i$ and
$\mu_i=\lambda_i/\lambda'$. For a test set $A\subseteq\NN^d$, use
Barbour's immigration--death solution, with its sign chosen so that
\[
 \mathcal A g(z)=\1_A(z)-\pi(A),\qquad
 \mathcal A g(z)=\sum_i\{\lambda_i\Delta_i g(z)
                           +z_i[g(z-\boldsymbol e_i)-g(z)]\}.
\]
At $z_i=0$ the corresponding death term is zero.
Write $H(z)=(\Delta_{ij}g(z))_{i,j}$ and
$c(\lambda')=(1+2\log^+(2\lambda'))/2$.
Barbour's Lemmas~2 and~3 \cite[p.~179]{Barbour1988SteinProcess}
give, for this same solution,
\[
 |\Delta_{ij}g(z)|\le1,\qquad
 |\alpha^T H(z)\alpha|
 \le c(\lambda')\sum_i\frac{\alpha_i^2}{\lambda_i}.
\]
The weighted bound, stated there for integer directions, extends to
rational directions by homogeneity and to real directions by continuity.
The literature input consists of the solution and these analytic bounds,
not an arithmetic or aggregated total-variation estimate.

For clarity, the constant entrywise bound also follows directly from the
solution's semigroup representation. Couple the processes started at
$z$, $z+\boldsymbol e_i$, $z+\boldsymbol e_j$ and
$z+\boldsymbol e_i+\boldsymbol e_j$ by adding two independent particles
of rate-one exponential lifetimes. The second difference vanishes once
either added particle dies; before that time the second difference of an
indicator is at most two in absolute value. Its time integral is thus
bounded by $\int_0^\infty 2e^{-2t}\,dt=1$. The argument includes $i=j$.

To obtain the weighted entrywise bound, conjugate the symmetric matrix
$H(z)$ by $D=\operatorname{diag}(\sqrt{\lambda_i})$.
The quadratic inequality gives $\|D H(z)D\|_{2\to2}\le c(\lambda')$.
Bounding each entry by this norm, and using the constant bound above,
yields
\begin{equation}\label{supp:eq:directional-factor}
 |\Delta_{ij}g(z)|\le
 \min\left\{1,\frac{1+2\log^+(2\lambda')}
                    {2\lambda'\sqrt{\mu_i\mu_j}}\right\}
 \le\min\left\{1,\frac{C(1+\log^+(2\lambda'))}
                    {\lambda'\sqrt{\mu_i\mu_j}}\right\}.
\end{equation}
The logarithmic term in multivariate Stein factors is discussed in
\cite{Rollin2012SteinFactors}; the two bounds used above are taken directly
from Barbour.
For the truncated geometric marks,
$S_E=\sum_{e=0}^Eq_e\in[1/2,1]$, $\lambda'=\lambda S_E$ and
$\mu_e=q_e/S_E$. The second expression is thus bounded by
$C(1+\log^+(2\lambda))/(\lambda\sqrt{q_eq_f})$.
The summability that removes dependence on the number of marks is
\begin{equation}\label{supp:eq:square-root-sum}
 \sum_{e\ge0}\sqrt{q_e}=1+\sqrt2.
\end{equation}

\begin{samepage}
\begin{theorem}[Finite aggregated comparison]\label{supp:thm:aggregated}
Let $X_e=\sum_{x\in I_N}K_{x,e}$, and let $D_Y(Q)$ be the bad maximal
supports. Put
\[
 s(\lambda)=1\wedge
          \frac{C(1+\log^+(2\lambda))}{\lambda}.
\]
Uniformly while $L+1$ and $Q+1$ remain in a fixed positive logarithmic
band and $2Q<Y$, one has
\begin{align}\label{supp:eq:aggregate-finite}
 &\EE\,\dTV\left(\Law((X_e)_{0\le e\le E}\mid\mathcal F_Y),
                  \bigotimes_{e=0}^E\Pois(\lambda q_e)\right)\notag\\
 &\hspace{1cm}\ll p\mathcal M_B+p|D_Y(Q)|
    +s(\lambda)p^2\bigl(NQ+E_Y(Q)+R_2(N,Q)\bigr).
\end{align}
Here $\mathcal M_B$ is the one-window defect mass at the base length,
and $E_Y(Q)$ is the ordered large-prime edge count for maximal supports.
\end{theorem}
\end{samepage}
\begin{proof}
Write $D=D_Y(Q)$ and retain the good-site counts
$G_e=\sum_{x\notin D}K_{x,e}$. Independently of the arithmetic variables,
introduce independent $Z_e\sim\Pois(\zeta_e)$ with
$\zeta_e=|D|pq_e$, and put $W=G+Z$. Removing the actual bad-site marks
and adding this Poisson filling changes the law by at most
\[
 \sum_{x\in D}\PP(J_{x,L}=1)+\sum_{e=0}^E\zeta_e
 \ll p\mathcal M_B+p|D|.
\]
This is an averaged conditional coupling bound. It has no factor $E$,
since the actual marks are disjoint refinements of $J_{x,L}$ and
$\sum_{e\le E}q_e\le1$.

Use the Stein solution for the \emph{full} target means $\lambda q_e$.
For the coordinate vector $\boldsymbol e_e$ and
$\Delta_eg(w)=g(w+\boldsymbol e_e)-g(w)$, Poisson summation gives,
conditionally on the good indicators and $\mathcal F_Y$,
\[
 \EE\bigl[Z_e\Delta_eg(W-\boldsymbol e_e)\mid G,\mathcal F_Y\bigr]
 =\zeta_e\EE\bigl[\Delta_eg(W)\mid G,\mathcal F_Y\bigr].
\]
Indeed $z\PP(Z_e=z)=\zeta_e\PP(Z_e=z-1)$ for $z\ge1$; the term with
$Z_e=0$ is zero. Thus the filling terms cancel the $\zeta_e$ part of the
birth rates in
\[
 \sum_{e=0}^E\{\lambda q_e\Delta_eg(W)
              -W_e\Delta_eg(W-\boldsymbol e_e)\}.
\]
Only the good indicators remain to be telescoped. The target parameters
are still $\lambda'=\lambda S_E$ and $\mu_e=q_e/S_E$, so
$\lambda'\sqrt{\mu_e\mu_f}=\lambda\sqrt{q_eq_f}$ exactly.
No comparison of the good-site intensity with $\lambda$ is needed,
even if every site was removed. The independent vector $Z$ stays in the
outside count during each telescoping step, so the good-site dependency
graph is unchanged. For good sites, apply the multivariate dependency-graph Stein argument
\cite{Barbour1988SteinProcess,Rollin2012SteinFactors} to the \emph{counts by mark}, not to
counts by site and mark. On the first-difference telescoping terms each
pair of types contributes its joint or product mass multiplied by
$|\Delta_{ef}g|$. The product masses are $p_ep_f$.
For local pairs use \eqref{supp:eq:local-marked}; for separated pairs use
\eqref{supp:eq:mixed-domination} after averaging. Decompose
$2^{\rho_Q}=1+(2^{\rho_Q}-1)$ and sum the second part by $R_2(N,Q)$.
The diagonal and all local pairs are bounded by $O(NQp^2)$.

Using the first bound in \eqref{supp:eq:directional-factor} sums the mark
weights to at most one. Using the second sums them to
$\sum_{e,f}\sqrt{q_eq_f}\le(1+\sqrt2)^2$. Take the better of these two estimates. This proves \eqref{supp:eq:aggregate-finite}
when $E\ge1$. For $E=0$, retain both signs during the same comparison,
so there are two positive target coordinates. At fixed excesses the
factor does not depend on the chosen signs; summing their probabilities
before the relative-pair bound again gives $R_2(N,Q)$.
The signed local factor is at most four, an absolute constant.
Thus the result also covers $E=0$ without extending the literature
statement to dimension one.
Notice that the arithmetic argument is used at $Q$, rather than merely
at $L$: after insertion of the directional weights, the unweighted
folding identity for marks is not the inequality being summed.
\end{proof}

At $Y=Y_N^\sharp$, the first two arithmetic error sources in
\eqref{supp:eq:aggregate-finite} are bounded by
\[
 \lambda e^{-V_N+o(\nu_N)}+N^{-1/2+o(1)}\lambda,
 \qquad
 \lambda(1+\log^+(2\lambda))e^{-V_N+o(\nu_N)}.
\]
For clarity, before absorbing subpolynomial factors the profile term is
\[
 p^2R_2(N,Q)
 \ll N^{-1/3+\varepsilon}
              \bigl(\lambda^2+\lambda 2^{E+2}\bigr).
\]
If $I+\log^+\lambda=O(V_N)$ and $E=O(V_N)$, multiplication by
$e^Is(\lambda)$ leaves only an $N^{o(1)}$ factor. Start with a smaller
arithmetic $\varepsilon$ to obtain the final remainder
$N^{-1/3+\varepsilon}$ at any prescribed positive $\varepsilon$.
This is the step that prevents an unrecorded exponential mark-cutoff loss.


For aggregation with signs retained, replace $q_e$ by
$q_{e,s}=q_e/2$ and $R_2(N,Q)$ by $R^{\rm val}_2(N,Q+1)$ in
\eqref{supp:eq:aggregate-finite}. The local factor is at most four and
$\sum_{e,s}\sqrt{q_{e,s}}=\sqrt2(1+\sqrt2)$, both fixed constants.
The full-value profile at $Q+1$ vertices has the same envelope as the
relative profile at $Q$ rows. It therefore leaves the same polynomial
remainder and the same one-factor budget. By contrast, the unweighted
labelled argument sums signs and marks first and continues to use
$R_2(N,L)$ at the base length.

\subsection{Removing the mark cutoff}
Assume $I+\log^+\lambda\le V_N-c\nu_N$, and choose
$0<c'<c''<c$. Set
\[
 E_N=\left\lceil
        \frac{\log^+\lambda+c''\nu_N}{\log2}\right\rceil .
\]
Then $E_N=O(V_N)=o(\log N)$, and the target tail has mass
$\mu_N=\lambda2^{-E_N-1}\le e^{-c''\nu_N}$.
The actual tail event is exactly a hit at length $L+E_N+1$.
The scalar hard-conditioning inequality at that shifted length gives
\[
 \PP(\exists x,\ e>E_N:K_{x,e}=1\mid C)
 \le \mu_N+
 O\left(e^I\mu_N
      (e^{-V_N+o(\nu_N)}+N^{-1/3+\varepsilon})\right).
\]
It is $O(\mu_N)$ under the displayed budget. Couple each tail to zero,
apply the finite theorem, and absorb
$\log(1+\log^+\lambda)=o(\nu_N)$. This gives
\[
 \dTV\left(\Law((X_e)_{e\ge0}\mid C),
                  \bigotimes_{e\ge0}\Pois(\lambda q_e)\right)
 \ll e^{-c'\nu_N}+N^{-1/3+\varepsilon}.
\]
For finite-mass sequences, tail summation is a measurable bijection,
with inverse $X_e=T_e-T_{e+1}$. The entire threshold path has exactly the
same total-variation distance.

\subsection{Keeping the spatial labels}
\label{supp:sec:information-cutoff}
For the labelled comparison no directional Stein factor is inserted.
The identities
\[
 \sum_e p_e\le p,\qquad
 \sum_{e,f}\PP(K_{x,e}K_{y,f}=1)\le\PP(J_{x,L}J_{y,L}=1)
\]
fold the separated-pair contribution at the base length. Their signed
versions hold after summing both signs. Thus the relation profile is
$R_2(N,L)$, whereas the graph and deleted set use the maximal support.

\paragraph{Free-cutoff bound and tails.}
Let $H=\log N$, $w\in[V_N,\widehat V_N]$, $Y=\lfloor e^w\rfloor$
and $\nu_w=H/w$. Suppose $C\in\mathcal F_Y$ has
$I=-\log\PP(C)\ge0$, and $\lambda=N2^{-L}\ge1$ with
$\ell=\log\lambda$ and $I+\ell=O(V_N)$. Choose
\[
 E=\left\lceil\frac{I+\ell+w}{\log2}\right\rceil,
 \qquad Q=L+E+1.
\]
Since $L=(H-\ell)/\log2$, the maximal number of values is
\[
 Q+1=\frac{H+I+w}{\log2}+O(1)=\log_2N+O(V_N).
\]
Both the base and maximal lengths remain in a fixed positive logarithmic
band, and $Y>2(Q+1)$ eventually, uniformly in these parameters.
By \cref{thm:marked-field} and the uniform free-cutoff estimates, the
finite field distance, multiplied by the information cost, is at most
\begin{equation}\label{supp:eq:information-free}
 C\left\{
 e^{I+\ell-D(H/w)+o(\nu_w)}+
 e^{I+2\ell-w+o(\nu_w)}+
 e^I\lambda(1+\lambda)N^{-1/3+\varepsilon/2+o(1)}
 \right\}.
\end{equation}
The weighted one-window term and the $N^{-1+o(1)}$ part of the edge
bound fit the polynomial term. Constants in the little-oh terms are
uniform for $w$ in the displayed interval and for the enlarged lengths.
A fixed multiplicative enlargement of this interval is also permitted by
\cref{supp:prop:cutoff-calculation}.

An excess greater than $E$ is exactly a threshold hit at $L+E+1$.
The unconditional first moment gives actual conditional tail cost at most
\[
 e^I\lambda2^{-E-1}(1+N^{-1/2+o(1)})=O(e^{-w}),
\]
and the target tail has no larger cost. This step uses conditioning only
as the bound $\PP(H\mid C)\le e^I\PP(H)$; no conditioning theorem at
a different cutoff is invoked. Coupling both tails to zero adds
$O(e^{-w})$ to \eqref{supp:eq:information-free}. The full site--sign--excess
space is countable, with finite expected total mass, so these tail
couplings exhaust its coordinates.

In the applications below, $I+\ell\le w\le\widehat V_N$. One may instead
use the common truncation $E=3\lceil\widehat V_N/\log2\rceil$, for which
$e^I\lambda2^{-E-1}\le \tfrac12 e^{-2\widehat V_N}$.
It still has $Q+1=\log_2N+O(V_N)$, so the finite comparison is unchanged
and both tails are $O(e^{-2\widehat V_N})$. This is the truncation used
in the proof of \cref{prop:information-labelled}; the sharper choice
above is not needed for that conclusion.

\paragraph{The information-dependent crossing.}
For $0\le I\le V_N$, set $F_I(w)=2D(H/w)-w-I$.
The function $D$ is continuous and strictly increasing on the larger
saddle branch. Since $H/w$ strictly decreases with $w$, $F_I$ is
continuous and strictly decreasing. Its endpoint values on
$[V_N,\widehat V_N]$ are $V_N-I$ and $-I$.
There is therefore exactly one root $w_N(I)$, including the endpoints
for $I=V_N$ and $I=0$. The first branch of
$g_{N,I}(w)=\min\{D(H/w)-I,(w-I)/2\}$ decreases strictly; the second
increases strictly. The root is their unique maximum on the interval.
It obeys
\[
 D(H/w_N(I))=\frac{w_N(I)+I}{2},\qquad
 \ell_{\rm lab,N}(I)=D(H/w_N(I))-I.
\]
If $\ell\le\ell_{\rm lab,N}(I)-c\nu_w$, the first two terms in
\eqref{supp:eq:information-free} are respectively
$e^{-c\nu_w+o(\nu_w)}$ and $e^{-2c\nu_w+o(\nu_w)}$.
The polynomial multiplier is $N^{o(1)}$ because
$I+2\ell\le w_N(I)=O(V_N)$. Also $w/\nu_w=w^2/H\asymp\log H\to\infty$,
so $e^{-w}$ is smaller than any fixed $e^{-c'\nu_w}$.
This proves \cref{prop:information-labelled}. The hypotheses are
vacuous at $I=V_N$ for $\lambda\ge1$, since then the allowed leading
logarithmic intensity is zero and the margin is strictly negative.

For $I>0$, the root satisfies $w_N(I)<\widehat V_N$. Hence
$2D(H/w_N(I))>2D(H/\widehat V_N)=\widehat V_N$, and its defining
equation gives $w_N(I)+I>\widehat V_N$.
Thus $w_N(I)>\widehat V_N-I$, giving the claimed strict improvement
of the leading budget over $\widehat V_N/2-I$ without differentiating
the implicitly defined root.
Uniformly for $w=tV_N$, $1\le t\le\widehat V_N/V_N$, the expansion
$D(\nu)=\nu u(\nu)(1+O(1/u(\nu)))$ gives
$D(H/w)/V_N=t^{-1}+o(1)$. The saddle equation consequently yields
$t^2+\theta t-2=0$ when $I/V_N\to\theta$.
The positive root is $(\sqrt{\theta^2+8}-\theta)/2$; subtracting
$\theta$ and dividing by two proves \eqref{eq:information-leading}.

\paragraph{A fixed event and a constrained cutoff.}
Suppose $C\in\mathcal F_{\lfloor e^{w_0}\rfloor}$ for a given
$w_0\in[V_N,\widehat V_N]$. The prime sigma-fields increase with the
cutoff. On the certified admissible interval $[w_0,\widehat V_N]$,
$g_{N,I}$ increases up to $w_N(I)$ and decreases afterwards; hence its
maximum is at $\max\{w_0,w_N(I)\}$. For any admissible $w$ the two
inequalities
\[
 I+\ell\le D(H/w)-cH/w,\qquad
 I+2\ell\le w-2cH/w
\]
are sufficient. In particular no event measurable only above the proposed
cutoff is retained by an algebraic optimization. All cutoffs here are
deterministic functions of the event and its probability, not of the
realized prime environment. This is an optimization of the proved
labelled bounds, not an optimality assertion about the law.

\subsection{Macroscopic transport and containment}\label{supp:sec:macro-marked}
Let $L_M=\lfloor\log_2M\rfloor-a_M$, $a_M=o(\log M)$, and
$\lambda=M2^{-L_M}$. The base-contained interior starts are
$2\le x\le M-L_M+1$; their number is $M-L_M$.
First retain $\lceil M^{1/2}\rceil\le x\le M-L_M+1$.
The conditional ranks and graph construction are local and unchanged.
The one-window defect masses, summed on dyadic slices, are
$M^{1/2+o(1)}$; the edge estimate uses the upper height $M$; and the
separated-pair mass is bounded by the macroscopic profile
\eqref{eq:master-envelope}. Thus the labelled finite comparison is
\[
 \ll\lambda(1+\lambda)
             \bigl(e^{-V_M+o(\nu_M)}+M^{-1/3+\varepsilon}\bigr).
\]
Restoring the deep starts costs the bound \eqref{eq:deep-starts},
\[
 O\left(\lambda M^{-1/2}
             +\exp(-c_0\log M/\log\log M)\right).
\]
The target deletion costs only $O(\lambda M^{-1/2})$.
Conditioning on a compatible event in the chosen prime field multiplies
the averaged error by at most $e^I$.

Choose $E=\lceil(I+V_M+\log(2+\lambda))/\log2\rceil$ in the hard
labelled regime. The exact mark-tail identity and the global scalar
bound control the real tail. Besides the intensity-linear tail term,
the scalar bound has the separate deep-start error just displayed;
that error is already present in the full-field inequality. In
particular it must not be silently dropped from a finite-prefix tail
calculation. Both tails are absorbed by
\eqref{eq:macro-field-rate}. The free-cutoff proof is identical with
$V_M$ replaced by the optimized cutoff in its two leading terms.

For aggregated counts apply \cref{supp:thm:aggregated} on the retained
macroscopic set, with $R_{2,\delta}(M,Q)$ in place of $R_2(N,Q)$.
The same geometric square-root sum applies. Under
$I+\log^+\lambda\le V_M-c\nu_M$, the deep-start correction times
$e^I$ is smaller than every fixed $e^{-c'\nu_M}$, and the aggregate
mark cutoff above removes both tails. This proves the macroscopic
one-factor staircase without an additional marked arithmetic input.

At exact centered level $r$ the target mean is precisely
$(M-L_M)2^{-b_M-r-1}$, not $M2^{-b_M-r-1}$.
All these fields retain the full rightward mark of a base-contained
start. They are not the list of blocks censored at the right end of the
finite sequence. The containment diagram in the article distinguishes
the two objects.

\subsection{Deterministic masks and diffuse limits}
A deterministic restriction of a lattice field contracts total
variation. At bounded intensity, if
$N^{-1}\sum_{x\in A_N}\delta_{x/N}\Rightarrow\mu$, the restricted marked
field therefore converges weakly to
$\operatorname{PPP}(\lambda\mu\otimes\nu_{\rm geo})$, along a subsequence
$N2^{-L}\to\lambda$. The same statement holds for the macroscopic field,
with scaling by $M$. It follows directly from the target Laplace
functional and weak convergence of the deterministic empirical measures.
There is no universal numerical rate for this last deterministic step
without a quantitative assumption on the mask or the test function.

Nor is the diffuse limit a total-variation target. The actual field is
supported on its lattice, while the diffuse target has positive
probability $1-e^{-\lambda\mu(\mathbb R)}$ of a nonempty configuration
outside that lattice. The approximation in total variation always takes
place \emph{before} passing from the grid to the continuum.
